Systematic uncertainties in the background estimates are driven by the uncertainties in the transfer factors, normalisations and shape of the tracklet \pt templates for the electron, muon, hadron and fake backgrounds. These uncertainties are dominated by the statistical uncertainties of the data in regions used to derive the background estimates, and are always below 10\%. When calculating values of parameters which are based on data, such as the TFs and smearing functions, the statistical uncertainty of the data is used as a $\pm 1\sigma$ variation. When performing the statistical fit, the values of the transfer factors and the individual bins of each tracklet \pt template are allowed to vary within their statistical uncertainties. Each source of background uncertainty (\pt template statistical uncertainties, TFs, and smearing functions) is treated as correlated across each set of associated CRs and SRs but uncorrelated between backgrounds, except $\mathrm{TF}^\mathrm{had}_{\mathrm{tracklet}}$, which is correlated with $\mathrm{TF}^e_{\mathrm{tracklet}}$ since both TFs are used in the estimation of the hadronic background. No dedicated uncertainty is associated with the overall background estimation methodology as, within other sources of uncertainties including the statistics of the data, the predicted post-fit tracklet \pt distributions in the VRs are consistent with the data.

Several sources of experimental uncertainty in the SUSY signal estimates are considered. The uncertainties are included in the likelihood fit as Gaussian nuisance parameters that affect the expectation values for each SR bin. Jet energy scale and resolution uncertainties are derived as a function of the jet $\pt$, $\eta$, and flavour, using data and simulated events, as detailed in Ref.~\cite{JETM-2018-05}.
Lepton reconstruction and identification uncertainties~\cite{EGAM-2021-01,MUON-2018-03} are included in the fit, but are found to have negligible impact.
Uncertainties in the tracklet reconstruction efficiency are estimated to be 10\% due to inactive pixel modules and 6\% due to selection differences between data and MC $V$+\;\!jets events. An uncertainty related to the low-energy charged pions is also applied, combining uncertainties of 10\% in the reconstruction efficiency, taken from Ref.~\cite{IDTR-2021-03}, and 6\% in the BDT selection efficiency, evaluated by varying the BDT input parameters by 1$\sigma$ and examining the effect on the output score. All uncertainties in the final-state object reconstruction are propagated to the $\met$ calculation, including an additional term taking into account uncertainties in the scale and resolution of the soft term~\cite{JETM-2020-03}. In the \SRAm region, the \met selection used is in the \enquote{turn-on} region of the \met trigger where the trigger is found to be between 85\% and 90\% efficient for the signal event topology. Therefore, a trigger efficiency uncertainty is also applied, as an uncertainty in the signal yield. 

The theoretical uncertainties originating from the choice of factorisation and renormalisation scales for the SUSY signal samples are evaluated by varying the corresponding MC generator parameters by a factor of two around their nominal values. The uncertainties associated with the choice of PDF set~\cite{Ball:2014uwa} and the value of the strong coupling constant are also considered. Finally, five variations of the A14 tune of \PYTHIA are used to account for ISR modelling effects due to the underlying event, jet structure, and different aspects of extra jet production. The total uncertainty from experimental sources is found to be ${\sim}5$\% for all signal scenarios. The theoretical modelling uncertainties contribute up to 20\%, increasing with increasing $\chinoonepm$/$\tau$-slepton mass, and are primarily driven by the ISR modelling uncertainty.

Figure~\ref{fig:systematicsUnc} presents a breakdown of systematic uncertainties affecting the background estimates in the SRs, grouped by the source of the uncertainty. The total uncertainty may not simply be the sum in quadrature of the individual uncertainties, because the fit can induce correlations between the components. In all SRs, the dominant uncertainty comes from the normalisation of the fake background, and the subdominant uncertainty in most regions arises from the shape of the fake background. Due to the limited number of tracklets available for the tracklet \pt templates in the \SRBo region, the statistical uncertainty of the template shape contributes significantly to the background uncertainty in this region. 

\begin{figure}[bt]
\centering
\includegraphics[width=0.55\textwidth]{figures/systematics/Draft_1/sr_syst.pdf}
\vspace{-0.75cm}
\caption{A summary of the systematic uncertainties affecting the background estimates in the SRs after the likelihood fit to data in the CRs is performed (\enquote{Post-fit}). The total uncertainty is indicated by a solid black line. The individual uncertainties may not sum in quadrature to the total uncertainty, due to correlations between the different components. The uncertainties arising from the estimation of the fake-tracklet background, which are the dominant uncertainties, are shown as blue dashed lines.}
\label{fig:systematicsUnc}
\end{figure}
